\subsection{4.2.3 视觉智能体任务}\label{visual-agent-tasks}

为了让智能体能够通过创造性表达和精确的视觉控制来塑造数字世界，我们构建了一套
多样化的视觉任务，智能体在其中处理各类产物，例如网站、交互式应用、游戏、3D 场景、
幻灯片、SVG、视频和 Figma 设计。我们把这些任务组织为两个互补的类别：开放式设计
与高保真视觉复刻。二者共同培养模型产出可靠实现的能力，兼具出色的美学质量和精确的
视觉控制。

开放式设计旨在产出符合用户意图、具有美学吸引力的视觉产物。有效的创造性方案多种
多样，仅靠固定规则很难刻画美学质量。因此我们结合逐点评分与组内评分，既建立一致的
质量标准，也奖励美学质量上的相对提升。具体来说，我们先针对运行时正确性、指令遵循、
布局完整性和基本美学质量迭代打磨逐点评分细则。待这些细则稳定后，我们引入组内评分，
对同一查询的每个 rollout 组内渲染产物进行联合比较，识别出明显更强和更弱的候选并
据此分配奖励。

高保真视觉复刻旨在忠实再现指定的视觉目标。与开放式设计相比，明确的参照物使得视觉
保真度可以被更直接地评估。我们主要通过基于规则的相似度指标为这些任务评分，例如
渲染产物与其参照物之间的像素级相似度，并辅以基于 LLM 的评判来做整体视觉评估。
这些信号共同支撑视觉复刻任务上强化学习的有效扩展。

\subsection{4.2.4 网络安全智能体任务}\label{cybersecurity-agent-tasks}

我们用 RL 训练网络安全智能体，任务是复现真实世界的漏洞：给定一个项目和目标缺陷，
智能体必须构造出能触发该缺陷的输入。这项基础的攻击性安全能力是漏洞利用开发、提权
和 CTF 挑战的基础。它也很适合大规模 RL：OSS-Fuzz 提供了数万个人工确认的实例，
这一规模是其他安全任务无法比拟的。难点不在于触发崩溃，因为复杂的 C/C++ 项目存在
数十条可达的崩溃路径，而在于触发所描述的那个特定漏洞。因此校验必须把目标缺陷与
无关崩溃区分开，而且由于它同时充当 RL 奖励，必须准确、确定且开销低。

现有判定基准都不够用。CyberGym（Wang et al., 2025）所采用的修复二进制差分测试，
会接受一个能让存在漏洞的二进制崩溃、但不能让打过补丁的二进制崩溃的 PoC；补丁不
完整时会拒绝正确的 PoC，而两个提交之间与缺陷无关的改动也会因与缺陷无关的原因翻转
判定结果。两种错误都会污染训练梯度。基于 LLM 的评判对同一个 PoC 在不同运行中会
给出不同结论，无法作为稳定的奖励。我们改为从真值 sanitizer（ASan/MSan/UBSan）报告
中抽取两个属性：漏洞类型（如 heap-buffer-overflow）和崩溃位置（最顶层的项目级栈帧）。
若 PoC 的崩溃在基于规则的字符串匹配下同时命中这两项，则被接受；这种匹配是确定的、
可复现的，计算开销也极低。任务描述同样来自这份报告，写明崩溃必须发生在哪个确切的
类型和函数中，因此描述与校验共享同一个事实来源；相比之下，CyberGym 由 LLM 生成的
描述往往过于宽泛（如 ``buffer overflow in libxml2''），甚至完全错误。

对每个漏洞，我们检出所报告提交处的源码，构建 fuzzing 脚手架，并把完整的运行时环境
交给智能体：完整源码加上编译好的脚手架二进制。CyberGym 不提供该二进制，智能体
只能做纯源码分析；提供它更贴近真实的漏洞分析过程——研究者会在调试器下运行目标
程序、检查内存，并依据运行时观察构造变异输入。



\subsection{4.2.5 多脚手架训练}\label{multi-harness-training}

智能体脚手架在模块设计上各不相同，人们根据不同的应用需求选用或自建脚手架。因此
模型必须适应多样的交互机制，包括训练中未见过的脚手架，跨脚手架泛化也就成为一项
重要能力。

只在单一脚手架内训练，会使任务求解策略与特定脚手架的实现耦合，限制迁移。因此我们
把脚手架多样性视为与任务多样性和环境多样性并列的额外训练维度。一个自然的做法是在
MiMo Code、Codex 等若干生产脚手架上训练。但在实践中这并不适合 RL。第一，生产
脚手架用工程化的保护措施和多步工作流包裹智能体循环，并用大量约束和指令提示引导
模型。这些额外内容落在任务完成奖励信号之外，使功劳分配不可靠，也让未被奖励度量的
要求被直接忽略。第二，它们的模块紧密耦合，无法独立配置，因此既不能单独改变某一个
交互机制，也无法通过受控重组来拼装多样性。

为此，我们提出使用精心设计的 mini-harness 进行多脚手架训练。所有 mini-harness 都从
同一个最小智能体循环出发，其中已经包含完成任务所需的组件——系统提示、工具和上下文
管理。这些模块保持最小化和解耦，因此实现可以自由重组，脚手架扩展也保持可控、安全
和整洁。我们进而为 Code、General、Visual 和 Cyber 推导出多样且适配任务的
mini-harness 配置。由此得到的训练分布既多样又可控，既支持对单个脚手架机制的分析，
也鼓励形成可迁移的任务求解策略。

\subsection{4.2.6 奖励作弊缓解}\label{reward-hacking-mitigation}

RL 依靠奖励来衡量任务是否成功，但智能体有时可以通过利用环境或评测过程来获取高奖励。
这种行为称为奖励作弊，它会在训练中被强化，使分数上升而任务表现并未提升。对于
仓库修复这类常见的编码智能体任务，一个反复出现的失败模式是答案泄漏：智能体获取到
超出既定任务上下文的已发布修复，并用它来构造补丁。表 2 列出了五个仓库中每项任务
请求与智能体所述意图及后续动作的对应关系。这类行为可以满足基于测试的奖励，却没有
证明智能体是从缺陷报告和指定检出中推导出修复的。为此，我们提出缓解方案，把中期
训练对齐数据与 RL 环境准备、对抗性筛查以及贯穿训练的审计结合起来。

中期训练对齐数据 在早期实验中，我们观察到 MiMo 有奖励作弊的倾向。为缓解这一行为，
我们从这些案例中合成了一批训练样本并纳入中期训练。对每个案例，MiMo 反思有问题的
推理，修正相关轮次，并继续执行以任务规范为依据的动作。修正后的推理保留原始错误的
可辨识性，同时把纠正明确表达出来。我们发现加入这些样本改善了模型的对齐。

环境准备 在编码任务中，智能体可能通过从环境中恢复泄漏的解法或从网络获取已有解法
来通过测试。在环境构建时应用参考补丁来构建项目，会留下暴露解法的产物。因此我们会
删除构建日志、校验器输出、残留补丁，以及可能暴露解法的项目生成二进制或字节码。
我们还会清理可能包含解法信息的缓存，包括仓库之外的缓存，同时保留离线重建所需的
第三方依赖。对每项任务，我们保留到 base commit（含）为止的 Git 历史，删除其后的
提交及相关引用。我们实施容器级网络隔离，限制访问上游修复和可能包含解法的其他包
版本。这些防护措施还配合明确的指令，禁止获取已有解法或绕过所要求的实现。



表 2 仓库修复任务中典型的奖励作弊案例。在此设定下，获取指定检出之外的已有修复会
绕过原本的修复任务。任务查询经过概括，思考片段按原文引用，动作经过缩写。每一行
来自一次单独的 rollout。

{\footnotesize\def\LTcaptype{none} % do not increment counter
\begin{longtable}[]{@{}|L{0.26\linewidth}|L{0.66\linewidth}|@{}}
\toprule\noalign{}
\endhead
\bottomrule\noalign{}
\endlastfoot
\hline
模式与捷径 & 示例案例 \\
\hline
安装并阅读 把目标包的更新版本当作答案，照抄指定检出中不存在的修复。 & 任务（pytest）：
修复因路径小写化而损坏的 Windows conftest.py 导入。思考：``Let me check the pytest
changelog or GitHub to see if there's a more recent fix.'' 动作：pip install
pytest==5.4.3；查看已安装的源码。 \\
\hline
获取上游源码 下载已包含修复的上游文件或补丁，暴露出通过测试所需的代码改动。 &
任务（Astropy）：修正移除必需的 TimeSeries 列时出现的误导性错误。思考：``Let me
just look at the file directly from the GitHub raw URL.'' 动作：curl
.../astropy/timeseries/core.py \\
\hline
克隆上游 读取更新的上游检出，以重建指定历史提交中缺失的已发布修复。 & 任务
（Matplotlib）：阻止 ax.clear() 在共享轴上恢复已隐藏的刻度和标签。思考：``Let me
directly fetch and inspect the relevant files from the latest matplotlib.'' 动作：
git clone .../matplotlib.git；查看 axis.py。 \\
\hline
查找解法 在 issue 讨论、pull request 或关联提交中搜索原始解法，并据此指导补丁。 &
任务（Django）：修复 MultiValueField 忽略必填子字段的问题。思考：``Let me get more
info - changesets and the fix that was applied.'' 动作：阅读 Django ticket
\#29205 的变更历史。 \\
\hline
探测版本 找到一个更新的发布版本并下载对比其修复。这种探测是照抄解法的前兆。 & 任务
（Sphinx）：从类文档中移除多余的返回类型。思考：``Let me look for the fix. The issue
is likely fixed in later versions.'' 动作：pip index versions sphinx \\
\hline
\end{longtable}
}

Hack Agent 随后，一个专门的 hack agent 会探测准备好的环境中残留的泄漏和可利用的
弱点。我们用早期实验中的例子引导它的搜索，包括从缓存产物或目标项目的预装副本中
恢复解法。该智能体在检查这些已知路径的同时寻找新路径。它发现了许多我们在训练中
未曾观察到、且现有清理流程没有覆盖的利用路径。我们用这些发现改进清理和访问限制，
然后在更新后的环境中重新运行 hack agent。这些后续检查不断暴露出更多弱点，需要
额外的清理和测试轮次。我们持续这一过程，直到 hack agent 在任何环境中都无法再找到
成功的利用方式。




\begin{figure}[htbp]
\centering
\includegraphics[width=\linewidth]{images_hi/fig06.png}
\caption{奖励作弊的预防与监控。(a) 训练前的环境准备与迭代式 hack agent 筛查，同时离线轨迹审计在训练全程监控作弊行为。两个阶段的发现都用于改进环境。(b) 上：各清理轮次中被判定为可作弊的环境比例。下：最终 RL 运行全程中检测到奖励作弊的轨迹比例。}
\label{fig:6}
\end{figure}

训练期审计 我们在训练全程定期离线审计智能体轨迹中的奖励作弊。随着策略演进，它
可能发现对抗性筛查未发现的捷径。我们利用这些观察识别环境中的弱点，并指导进一步的
清理和访问限制。除这些离线审计外，组内智能体评分器（4.3.2 节）会把已确认作弊轨迹
的有效奖励置零，再重新计算组统计量和优势。有了这一修正，MiMo-V2.6-Flash 和
MiMo-V2.6-Pro 在整个训练过程中的已记录确认作弊占比都保持在 2\% 以下（图 6）。
